\section{Limitations}
\label{sec:limitations}

Several limitations should be considered before treating the reported numbers as universal. First, the study is site-specific. The plant is located in a tropical coastal setting, and its cloud, humidity, temperature, and monsoon-season behavior may differ from inland, desert, or high-latitude PV sites. Transfer to another plant would require retraining or at least calibration.

Second, the model depends on the quality of the weather inputs. The training dataset contains aligned variables, but deployed forecasts will often come from numerical weather prediction or API-based weather services. Forecast cloud cover can be spatially coarse and may miss local cloud-edge timing. This matters because most large residuals occur during ramp-heavy intervals.

This weather-input dependency is also why the sensor-assisted diagnostic result should not be advertised as the operational result. The sensor model answers an engineering question: if target-time irradiance is known, how well can the plant conversion behavior be fitted? The proposed forecast path answers a different question: if only weather and cloud inputs are available for the target timestamp, how well can power be estimated? These two settings must remain separate in the abstract, tables, and conclusion.

Third, the reported model is trained and evaluated on rows retained after data cleaning. Machine-fault periods, hardware-level errors, plant shutdown intervals, and invalid sensor states were omitted so that the learning target represents normal machine operation and weather-driven PV variability. If future datasets include unmarked downtime, inverter clipping faults, data-acquisition gaps, or corrupted sensor periods, those intervals should again be filtered, flagged, or modeled separately. Otherwise, the model may learn equipment failures as if they were atmospheric losses.

The data-quality table verifies cadence, duplicates, missing raw values, basic range behavior, and the omission of machine/hardware fault periods during cleaning. It does not replace a full SCADA-event audit. For example, a plant can have plausible non-missing numeric values during a curtailment, maintenance event, or inverter derate. Such periods require event logs, inverter status, or rule-based operating-state filters when new data are added.

Fourth, the sensor-feature path obtains especially strong metrics because it has access to closely aligned operational measurements. Such results help estimate the achievable plant-data fit, but day-ahead deployment should emphasize the forecast-input path. This separation should remain clear when interpreting the reported results.

Fifth, the solar-time implementation is simplified. The feature pipeline uses IST civil time directly with a declination and hour-angle approximation; it does not apply the equation of time or a longitude correction. This is reproducible and adequate for the reported model, but later work should compare it with a pvlib-style solar-position implementation.

Finally, the ordered references emphasize microgrid operation, building energy management, demand response, and residential energy scheduling. They motivate the operational value of PV forecasts, but they do not by themselves validate every modeling choice in the forecasting pipeline.
